%!TeX program = pdflatex
\documentclass{resume}

\usepackage{tikz, lmodern, amsmath, amssymb}

\definecolor{papercolor}{HTML}{282828}
\definecolor{textcolor}{HTML}{eeeeee}

%\usepackage[pagecolor={papercolor}]{pagecolor}
%\color{textcolor}


\usepackage{fontawesome5}

\newcommand{\bars}[1]{
    \tikz[overlay, remember picture] \foreach\i in {1,...,3}
        \draw ({(\i - 2)*4pt},0) rectangle +(2pt,{(\i + 1)*2 pt});
    \tikz[overlay, remember picture] \foreach\i in {1,...,#1}
        \fill ({(\i - 2)*4pt - 2.5pt},0) rectangle +(2pt,{(\i + 1)*2 pt});
}

\begin{document}

\fontfamily{ppl}\selectfont

\noindent
\begin{tabularx}{\linewidth}{@{}m{1\textwidth}}%{@{}m{0.78\textwidth} m{0.2\textwidth}@{}}
{
    \Large{Isabella Basso do Amaral} \newline
    \small{
        \clink{
            \href{mailto:isabellabdoamaral@usp.br}{\faIcon{envelope}\ isabellabdoamaral@usp.br}
            \href{https://telegram.me/isinyaaa}{\faIcon{telegram}}
            \href{https://gitlab.com/isinyaaa}{\faIcon{gitlab}}
            \href{https://github.com/isinyaaa}{\faIcon{github} @isinyaaa}
        } \newline
          São Paulo, Brazil\newline
        {\footnotesize Linux \faIcon{linux} enthusiast \textbf{·} \faIcon{heart} FOSS}
        %{\footnotesize *Rio de Janeiro durante a quarentena}
    }
} %&
%{
%    \centering
%    \frame{\includegraphics[width=3cm]{images/photo.png}}
%}
\end{tabularx}
\begin{center}
\begin{tabularx}{\linewidth}{@{}*{2}{X}@{}}
% left side %
{
    \csection{Work experience}{\small
        \begin{itemize}
            \item \frcontent{Software Engineer}
            {Red Hat}
            {Django, SQL, Python, Containers, Bash, C, Git, Remote team collaboration, CI}
            {
                I currently work with static analysis CI solutions for Fedora/RHEL RPM packages.
            }
            {November 2021 onwards}
            \item \frcontent{Backend Developer}
            {Data Machina}
            {Python, Google Cloud, Firebase, Docker, Git, Agile Methods}
            {
                - I mentored an intern;

                - led the development of an AR solution;

                - previously was developing a WhatsApp chatbot;

                Team leading talent;

                Assisted and led various key maintenance tasks (deploying, setting up CI, etc.).
            }
            {February 2021 up to October 2021}
            \item \frcontent{Programmer}
            {Freelancer}
            {Python \& Java}
            {
                - Part-time computer science tutoring;

                - Freelancing for scientific applications.
            }
            {July 2020 up to January 2021}
        \end{itemize}
    }
    \csection{Side projects}{\small
        \begin{itemize}
            \item \frsimple{LKCAMP group \clink{\href{https://lkcamp.dev/}{[lkcamp.dev/]}}}
            {C, Linux}
            {Group member, independent Linux Kernel contributor}
            \item \frsimple{Kernel mentorship}
            {Bash, C, Linux}
            {Mentored by AMD GPU engineer, independent Linux Kernel contributor}
            %\item \frsimple{Data Girls Neuron\clink{\href{https://www.linkedin.com/company/data-girls-neuron/}{[Linked In]}}}
            %{Data Science, events and community support}
            %{Member}
            \item \frsimple{INOVA USP lab}
            {Bash, Linux}
            {Sys-admin}
            \item \frsimple{Cryptography and Cybersecurity study club}
            {Cryptography, Cybersecurity}
            {Group member}
            %{
            %    Busco me qualificar como sys-admin por utilizar Linux no dia-a-dia e
            %    me interessar pelo funcionamento do sistema, assim como em utilizá-lo
            %    plenamente
            %}
            %{}
            %\item \frsimple{Technical illustration design}
            %{Art}
            %{\LaTeX}
            %{
            %    Produzi {\href{https://github.com/isinyaaa/werk-yago-pub}{material didático}} para diversas disciplinas
            %    de curso técnico;
            %
            %    Atualmente faço ilustrações técnicas para um
            %    projeto de {\href{https://petfisica.home.blog/olimpet/}{ensino de física}}
            %}
            %{}
            %\item \frsimple{Translation}
            %{English}
            %{english--portuguese}
            %{
            %    Traduzi diversas obras de
            %    filosofia para publicação (artigos e livro)
            %}
            %{}
            %\item \frsimple
            %{Algorithmic math problem resolutions \clink{\href{https://projecteuler.net/}{[projecteuler.net/]}}}
            %{Maths}
            %{Julia, Haskell \& Python}
            %{
            %    Resolvi diversos problemas do acervo para afiar minhas
            %    habilidades analíticas
            %}
            %{}
            %\item \frsimple{Bare Windows' UIs \clink{\href{https://github.com/isinyaaa}{[github.com/isinyaaa]}}}
            %{Tinkering}
            %{C\# \& XAML}
            %{
            %    Fruto de um interesse de longa data em interfaces e interações para
            %    além do CLI
            %}
            %{}
        \end{itemize}
    }
}
% end left side %
&
% right side %
{
    \csection{Education}{\small
        \begin{itemize}
            % item 1 %
            \item \frcontent{Molecular Sciences BSc\newline emphasis on maths \& computer science}
            {Universidade de São Paulo}
            {Ongoing}
            {
                As a scientist by nature, I thrive in interdisciplinary contexts;
                \href{http://cecm.usp.br/}{that's exactly what my degree is about}.

                Scientific guidance student:

                - currently looking into the \href{https://isinyaaa.github.io}{Linux Kernel \& its graphics stack};
                %- currently looking into FPGAs \& optimization;

                - previous topics were algorithms \& \href{https://github.com/isinyaaa/uni-latex/tree/main/combinatorics}{graph theory}.
            }
            {2020 onwards}
            %\item \frcontent{Associate degreee in Data Science}
            %{Universidade Estácio de Sá}
            %{Ongoing}
            %{
            %    I started studying data science topics to sharpen my
            %    understanding and improve my skills.
            %}
            %{from 2021 onwards}
            %\item \frcontent{Physics BSc}
            %{Universidade de São Paulo}
            %{Interrupted}
            %{
            %    The degree I started at;

            %    As a passionate student I used to wonder about the HOWs;

            %    Then I found a degree more in line with my goals.
            %}
            %{2020}
            \item \frcontent{Industrial Automation Technical School}
            {Instituto Federal Fluminense}
            {Concluded}
            {
                Organized \href{http://lattes.cnpq.br/9507659630401803}{several projects and scientific presentations};

                I was my class' top student
            }
            {2017 -- 2019}
        \end{itemize}
    }
    \csection{Skills}{\small
      	\begin{itemize}
      		\item \textbf{Programming}
      		{\footnotesize
      			\begin{itemize}
      			\item[{\bars{3}}] ($\geqslant$2y): \LaTeX\, (3y) \& Python (4y)
      			\item[{\bars{2}}] ($\geqslant$1y): Bash, C, C\# \& Java
      			\item[{\bars{1}}] ($<$1y): C++, Haskell \& Rust
      		\end{itemize}}
      		\item \textbf{Languages}
      		{\footnotesize
      			\begin{itemize}
      			\item[{\bars{3}}] Portuguese (native), English (C1 + 6mo interchange + translation experience)
      			\item[{\bars{2}}] Spanish
      			\item[{\bars{1}}] French \& German
      		\end{itemize}}
      	\end{itemize}
    }
}
\end{tabularx}
\end{center}
\end{document}